\uuid{4Esn}
\titre{Ensembles}
\auteur{bodin ; croizat ; sacre}
\organisation{exo7}
\licence{CC BY-NC-SA 4.0}
\datecreate{2022-03-08}
\niveau{L1}
\module{Algèbre}
\chapitre{Logique, ensemble, raisonnement}
\sousChapitre{Ensemble}
\difficulte{1}
\theme{Ensembles}
\exo7qcmid{hilisit.logique.023}

\contenu{
\begin{question}
\qclassification{100.02}
Soit la fonction $f : \Rr \to \Rr$ définie par $f(x) = x^2-2$.
Quelles sont les affirmations vraies ?
\begin{answers}
	\bad{L'image de $-2$ est $-6$.}
	\good{Un antécédent de $7$ est $-3$.}
	\bad{La valeur $-2$ admet plusieurs images.}
	\good{La valeur $7$ admet plusieurs antécédents.}
\end{answers}
\begin{explanations}
L'image de $-2$ n'est pas $-6$ car $f(-2) = (-2)^2-2=4-2=2$. \\
Le nombre $-3$ est antécédent de $7$ car $f(-3) = 7$. \\
La valeur $x=-2$ admet une seule image. C'est en fait vrai pour tout $x\in\Rr$ !\\
Les antécédents de $7$ sont les solutions de l'équations $f(x) = 7$, c'est-à-dire $x^2-2=7$ ou encore $x^2=9$, ce sont donc $x=3$ et $x=-3$. Il y a donc deux antécédents à la valeur $7$.
\end{explanations}
\end{question}
}
